\section{Related work}
\label{sec:related}

\subsection{Concept erasure}
Concept erasure aims to suppress designated concepts in a pretrained text-to-image model while preserving its ability to generate unrelated content.
Early approaches achieve this by updating model parameters.
ESD~\cite{Gandikota_2023_ICCV} fine-tunes diffusion weights using a target derived from negative guidance to reduce the generation of the erased concept.
UCE~\cite{Gandikota_2024_WACV} computes closed-form updates to cross-attention projections, enabling multi-concept editing.
Concept Ablation~\cite{Kumari_2023_ICCV} fine-tunes the model to align the conditional distribution of a target concept with that of an anchor concept.
Extending weight editing to rectified flow transformers, EraseAnything~\cite{pmlr-v267-gao25j} combines LoRA-based bi-level optimization with attention regularization and self-contrastive learning to balance erasure and preservation.
In contrast, inference-time methods avoid modifying model weights: SLD~\cite{Schramowski_2023_CVPR} keeps parameters fixed and introduces an additional safety-guidance term during sampling to suppress inappropriate content.

Our method builds upon this inference-time paradigm, encoding suppression within a learned negative condition while preserving the original model weights.
Unlike indiscriminate guidance, we pair each suppression prompt with a separately learned detection prompt that selectively determines when to activate the intervention.

\subsection{Negative prompts}
Negative prompting specifies content to avoid through the negative condition in classifier-free guidance (CFG), steering generation away from the corresponding prediction.
Beyond manually written negative text, recent work explores learned conditions and modified guidance rules.
Shin \etal~\cite{shin2025promptbasedsafetyguidanceineffective} employ concept inversion~\cite{pham2023circumventing} to recover implicit negative embeddings, restoring safety guidance in models whose responses to explicit concept prompts have been weakened by unlearning.
SafeCFG~\cite{pan2025safecfg} replaces the null-text embedding in CFG with an input-dependent embedding generated by a learned Transformer to suppress harmful content.
ReNeg~\cite{Li_2025_CVPR_ReNeg} optimizes negative embeddings via reward guidance to improve image quality and human preference.
TraSCE~\cite{jain2026rethinkingnegativepromptingconcept} combines unconditional anchoring with alignment-localized steering to counter guidance collapse, where positive and negative predictions converge.

These studies motivate learning effective negative representations and controlling how they enter sampling.
Rather than predicting embeddings per input or optimizing for general aesthetic rewards, we learn a dedicated suppression prompt per concept and reuse it across inputs flagged by the detection prompt, directly optimizing it for concept suppression within a frozen generative model.

\subsection{Prompt tuning}
Prompt tuning adapts pretrained models through learned prompt representations while keeping their backbones frozen.
Lester \etal~\cite{lester-etal-2021-power} learn soft prompt embeddings prepended to language-model inputs, enabling task adaptation with a small set of trainable parameters.
CoOp~\cite{zhou2022coop} replaces hand-crafted context words with learnable vectors to adapt frozen CLIP models to downstream image recognition tasks.
Textual Inversion~\cite{Gal_TextualInversion} learns new token embeddings from a few reference images to represent specific objects or styles in a frozen text-to-image model.
PEZ~\cite{NEURIPS2023_a0054803} combines gradient-based optimization with vocabulary projection to discover interpretable discrete prompts for text and image generation.

Following this parameter-efficient paradigm, we apply continuous prompt learning to selective concept suppression while keeping both the text encoders and the generative model entirely frozen.
Specifically, we decouple the learned representations into a detection prompt for input discrimination and a suppression prompt for negative conditioning during generation, optimizing each independently for its respective role.